% =====================================================================
%  Chapitre 6 — Sprint 4 : Abonnements, intégrations et mise en production
% =====================================================================
\chapter{Sprint~4 : Abonnements, intégrations et mise en production}
\label{chap:sprint4}

\section*{Introduction}
Ce dernier sprint transforme une plateforme fonctionnellement complète en un
service exploitable et commercialisable. Il met en place le système
d'abonnements et de facturation, l'application des quotas par plan, la
supervision par le super-administrateur, la couche d'intégration permettant
d'embarquer l'éditeur, le site marketing public, et enfin le déploiement en
production. Nous présentons les objectifs du sprint, les cas d'utilisation,
les diagrammes de séquence, puis la réalisation.

\section{Objectifs du sprint}
Ce sprint couvre les epics~11 à~15 du backlog. Le
tableau~\ref{tab:backlog-s4} en présente les user stories.

\begin{table}[H]
\centering
\caption{Backlog du sprint~4}
\label{tab:backlog-s4}
\small
\begin{tabularx}{\linewidth}{clX}
\toprule
\textbf{ID} & \textbf{User story} & \textbf{Description} \\
\midrule
11.1 & Souscrire un abonnement & En tant qu'administrateur, je veux souscrire un plan payant via un paiement en ligne sécurisé. \\
11.2 & Gérer l'abonnement & En tant qu'administrateur, je veux changer, annuler ou réactiver mon abonnement. \\
11.3 & Quotas par plan & En tant que système, je veux appliquer les limites propres à chaque plan. \\
12.1 & Attribuer les plans & En tant que super-administrateur, je veux superviser les entreprises et leur attribuer un plan. \\
13.1 & Clés d'API & En tant qu'administrateur, je veux générer des clés d'intégration. \\
13.2 & Embarquer l'éditeur & En tant qu'outil externe, je veux envoyer mes utilisateurs dans l'éditeur sans seconde authentification. \\
14.1 & Site marketing & En tant que visiteur, je veux consulter un site public présentant la plateforme et ses offres. \\
15.1 & Mise en production & En tant qu'équipe, je veux déployer la plateforme de façon sécurisée et reproductible. \\
\bottomrule
\end{tabularx}
\end{table}

\section{Raffinement et description des cas d'utilisation}

\subsection{Cas d'utilisation « Gérer son abonnement »}
La figure~\ref{fig:uc-billing} regroupe les cas d'utilisation liés à la
facturation : souscription, changement de plan, annulation, réactivation, et
attribution de plan par le super-administrateur.

% [FIGURE : diagramme de cas d'utilisation « Gérer son abonnement »]
\figplaceholder{Cas d'utilisation « Gérer son abonnement »}{fig:uc-billing}{7cm}

Le tableau~\ref{tab:uc-subscribe} décrit le scénario « Souscrire un
abonnement ».

\begin{table}[H]
\centering
\caption{Description textuelle du cas d'utilisation « Souscrire un abonnement »}
\label{tab:uc-subscribe}
\small
\begin{tabularx}{\linewidth}{lX}
\toprule
\textbf{Cas d'utilisation} & Souscrire un abonnement \\
\textbf{Acteur} & Administrateur d'entreprise \\
\textbf{Préconditions} & L'administrateur est connecté~; son entreprise est sur un plan inférieur. \\
\textbf{Postconditions} & L'entreprise est passée au plan payant choisi. \\
\midrule
\textbf{Scénario nominal} &
1.~L'administrateur choisit un plan et une périodicité (mensuelle ou annuelle).\newline
2.~Le système ouvre le paiement intégré et présente le détail hors taxes et TTC.\newline
3.~L'administrateur renseigne son moyen de paiement et confirme.\newline
4.~Le prestataire de paiement valide la transaction.\newline
5.~Le système met à jour le plan de l'entreprise et débloque les fonctionnalités correspondantes. \\
\midrule
\textbf{Scénario alternatif} &
4.a.~En cas d'échec du paiement, le système en informe l'administrateur, qui peut réessayer. \\
\bottomrule
\end{tabularx}
\end{table}

\section{Diagrammes de séquence}

\subsection{Séquence « Souscrire un abonnement »}
La figure~\ref{fig:seq-subscribe} décrit l'enchaînement entre le client, la
passerelle API, le prestataire de paiement et le service
d'authentification lors d'une souscription, jusqu'à la mise à jour du plan
de l'entreprise.

% [FIGURE : diagramme de séquence « Souscrire un abonnement »]
\figplaceholder{Diagramme de séquence « Souscrire un abonnement »}{fig:seq-subscribe}{7cm}

\subsection{Séquence « Embarquer l'éditeur »}
La figure~\ref{fig:seq-embed} illustre le flux d'intégration : un outil
externe redirige son utilisateur vers l'éditeur avec une authentification
déléguée, celui-ci construit un modèle, puis est redirigé vers l'outil
d'origine.

% [FIGURE : diagramme de séquence « Embarquer l'éditeur »]
\figplaceholder{Diagramme de séquence « Embarquer l'éditeur »}{fig:seq-embed}{7cm}

\section{Réalisation}

\subsection{Plans et grille tarifaire}
Quatre plans ont été définis : \textbf{Free} (gratuit), \textbf{Pro},
\textbf{Pro Organisation}, ainsi qu'un plan \textbf{interne} non
commercialisable réservé aux entités du groupe. Le
tableau~\ref{tab:plans} récapitule leurs caractéristiques.

\begin{table}[H]
\centering
\caption{Plans d'abonnement de la plateforme}
\label{tab:plans}
\small
\begin{tabularx}{\linewidth}{lcX}
\toprule
\textbf{Plan} & \textbf{Tarif (HT/mois)} & \textbf{Fonctionnalités} \\
\midrule
Free & Gratuit & Canal e-mail uniquement, 1 modèle, 1 interaction IA. \\
Pro & 25~€ & Tous les canaux, IA illimitée (hors invitations et clés d'API). \\
Pro Organisation & 55~€ & Pro + gestion d'équipe + accès API / intégrations. \\
Interne & — & Illimité~; attribué par le super-administrateur aux entités du groupe. \\
\bottomrule
\end{tabularx}
\end{table}

La grille tarifaire publique (figure~\ref{fig:ui-pricing}) propose une
bascule entre facturation mensuelle et annuelle. La \textbf{TVA} française
(20~\%) est correctement appliquée, le détail hors taxes et toutes taxes
comprises étant présenté au moment du paiement.

% [FIGURE : capture — page de tarification]
\figplaceholder{Page de tarification}{fig:ui-pricing}{6cm}

\subsection{Paiement et facturation}
Le paiement s'appuie sur le prestataire Stripe, au moyen d'un
\textbf{Checkout intégré} : le règlement s'effectue sur la page de la
plateforme, sans redirection externe (figure~\ref{fig:ui-checkout}). Les
abonnements sont récurrents et peuvent être mensuels (flexibles) ou annuels
(engagement de douze mois, à tarif réduit). Un onglet de facturation permet à
l'administrateur de consulter son plan, de le changer, de l'annuler dans le
respect de l'engagement, ou de le réactiver. Les événements de paiement sont
répercutés sur le plan de l'entreprise, tant par confirmation directe que par
\emph{webhook} de sécurité.

% [FIGURE : capture — paiement intégré Stripe et onglet facturation]
\figplaceholder{Paiement intégré et onglet de facturation}{fig:ui-checkout}{6cm}

\subsection{Application des quotas par plan}
Les limites de chaque plan sont appliquées côté passerelle API à partir
d'une source unique de vérité. Le plan gratuit est restreint au canal
e-mail, à un modèle et à une interaction IA, ce décompte étant monotone (il
n'est jamais décrémenté, de sorte qu'une suppression puis recréation ne
permet pas de contourner la limite). Lorsqu'une limite est atteinte, une
\textbf{fenêtre de montée en gamme} invite l'utilisateur à passer à un plan
supérieur (figure~\ref{fig:ui-upgrade}).

% [FIGURE : capture — fenêtre de montée en gamme]
\figplaceholder{Fenêtre de montée en gamme}{fig:ui-upgrade}{5cm}

\subsection{Supervision par le super-administrateur}
L'espace de super-administration a été enrichi d'une vue des commandes et
d'indicateurs clés (revenu récurrent, entreprises actives, impayés) obtenus
par enrichissement des données de facturation. Le super-administrateur peut y
attribuer un plan à toute entreprise, y compris le plan interne. La
figure~\ref{fig:ui-superadmin-billing} présente cet espace.

% [FIGURE : capture — supervision des abonnements (super-admin)]
\figplaceholder{Supervision des abonnements}{fig:ui-superadmin-billing}{6cm}

\subsection{Intégrations et embarquement}
Un administrateur peut générer des \textbf{clés d'API} depuis le panneau
« Intégrations » des paramètres. Un \textbf{flux d'embarquement} permet à un
outil externe (par exemple un CRM) d'envoyer ses utilisateurs dans l'éditeur
avec une authentification déléguée — sur le principe d'un Checkout délégué,
sans seconde connexion —, de leur faire construire un modèle, puis de les
rediriger vers l'outil d'origine, le tout cloisonné à l'entreprise
concernée.

\subsection{Site marketing public}
Un site marketing public a été développé pour présenter la plateforme
(accueil, fonctionnalités, tarifs, à propos, contact, démonstration). Sa page
de tarification est reliée aux plans réels de la plateforme. La
figure~\ref{fig:ui-marketing} en présente la page d'accueil.

% [FIGURE : capture — site marketing (accueil)]
\figplaceholder{Site marketing public}{fig:ui-marketing}{6cm}

\subsection{Mise en production}
L'ensemble des services est \textbf{conteneurisé} avec Docker et orchestré
par Docker Compose, deux piles distinctes étant maintenues pour le
développement et la production. La plateforme est \textbf{déployée en
production} sur un serveur du groupe. Son exposition publique repose sur un
\textbf{tunnel Cloudflare}, sécurisé et sortant, qui ne nécessite l'ouverture
d'aucun port entrant et confie la terminaison HTTPS à Cloudflare. Plusieurs
noms d'hôte publics desservent respectivement le frontend, la passerelle API,
le stockage MinIO et les services d'IA. Le code est hébergé sur GitLab, les
secrets étant tenus hors du dépôt. La figure~\ref{fig:archi-deploy} présente
l'architecture de déploiement.

% [FIGURE : architecture de déploiement (Docker + Cloudflare Tunnel)]
\figplaceholder{Architecture de déploiement}{fig:archi-deploy}{6.5cm}

\section*{Conclusion}
Ce sprint a doté la plateforme des attributs d'un véritable service SaaS de
production : un système d'abonnements Stripe avec TVA et gestion complète du
cycle de vie, l'application rigoureuse des quotas par plan, la supervision
par le super-administrateur, une couche d'intégration et d'embarquement, un
site marketing public, et un déploiement sécurisé via Docker et Cloudflare.
La plateforme est ainsi passée de la preuve de concept fonctionnelle au
produit exploitable en conditions réelles.
